\chapter{Research Objectives and Hypotheses}
\label{sec:research-objectives}

\section{Research Problem}
\label{sec:research-problem}

The central question of this thesis is not whether neural relays are
universally superior to classical relay strategies, but under what conditions
they can improve upon or complement them. The comparison depends on more than
bit-error rate (\texttt{BER}) alone. A classical detector may have an
appropriate model-based decision rule but require channel information that is
unavailable or expensive to estimate. Conversely, the required information
may be available while the corresponding computation or decision delay becomes
prohibitive.

The research problem is therefore organized around two distinct boundaries:
\emph{identifiability} and \emph{arithmetic}. Identifiability concerns whether
the information required to construct the appropriate model-based classical
rule is supplied or can be estimated with sufficient reliability. Arithmetic
concerns whether that rule, once constructible, can be executed within the
available computational and latency budget.

These boundaries need not occur at the same operating point. A classical
detector may remain statistically appropriate while becoming computationally
expensive as channel memory or modulation order increases. In another regime,
the computation may remain manageable while the channel information required
to instantiate the detector cannot be estimated reliably from the available
observations.

The thesis therefore asks two related questions:

\begin{enumerate}
    \item Under what information conditions does learned relay processing
    improve upon or complement the strongest applicable classical baselines
    evaluated in this thesis?

    \item When a model-based classical solution remains constructible, under
    what computational or decision-delay conditions does its implementation
    become operationally restrictive?
\end{enumerate}

The objective is consequently not to identify a universally winning relay
architecture. It is to determine the regimes in which the requirements of
classical and learned processing remain satisfied.


\section{Research Objectives}
\label{sec:research-objectives-detail}

The experimental program is organized around three objectives.

\paragraph{O1: Establish the matched-model control.}

Determine whether learned relay processing provides a measurable BER advantage
when the assumptions required by classical symbol-wise processing are fully
satisfied. The canonical benchmark therefore fixes the source, channel model,
modulation, destination processing, and evaluation procedure while varying the
relay processing function.

This objective establishes the reference regime against which the later
departures are interpreted.

\paragraph{O2: Locate the identifiability boundary.}

Determine how the comparison between learned and classical processing changes
as the information required by model-based detection becomes progressively
more difficult to obtain. This objective is evaluated by first introducing
channel memory while retaining exact channel knowledge, then replacing exact
knowledge with finite pilot-based estimation, and finally removing the
supplied per-block information required to construct the matched posterior.

The purpose is not to weaken the classical algorithms themselves. At each
stage, the comparison uses the strongest applicable classical baseline
evaluated under the information available in that regime.

\paragraph{O3: Locate the arithmetic boundary.}

Determine when an available model-based solution becomes operationally
expensive under computational and decision-delay constraints. This objective
compares the scaling and latency requirements of sequence-based classical
processing, learned inference, and block-based relay processing where those
constraints are relevant.

The arithmetic objective is distinct from the identifiability objective. A
classical rule may be constructible but expensive, or inexpensive but
impossible to instantiate reliably because the required channel information
is unavailable.


\section{Experimental Regime Map}
\label{sec:experimental-regime-map}

The experiments are organized as an ordered regime map rather than a strict
one-variable-at-a-time ablation. The sequence progressively relaxes the
preconditions required by model-based classical processing.

The underlying two-hop relay problem remains the organizing communication
system, while the channel model, modulation, estimator, evaluation budget, and
strongest applicable classical baseline are adapted where required by the
question posed in each regime. The sequence should therefore be interpreted
as a progression in the information available to the receiver rather than as
a claim that every other simulator parameter remains unchanged.


\paragraph{Layer 1: matched, memoryless control.}

Layer~1 is the canonical matched-model control. The channel is memoryless and
perfectly known, so the classical symbol-wise decision rule has the model
information required by the problem.

This layer asks whether learning provides an advantage when the preconditions
of classical processing are fully satisfied. A failure of the learned relay
to outperform the matched classical rule in this regime is therefore
informative rather than evidence that learned relaying is generally
ineffective. It establishes the reference point from which the later
departures are interpreted.


\paragraph{Layer 2: matched channel with memory.}

Layer~2 introduces channel memory while retaining exact channel knowledge.
The required classical decision therefore changes from a symbol-wise rule to
a sequence-dependent rule.

This layer primarily tests the effect of function class and computational
structure. The channel is not unknown at this stage. Both the learned and
classical methods are evaluated with the information required by the
corresponding experiment, while the classical reference becomes a
sequence-based detector appropriate to the known finite-memory channel.

Layer~2 therefore separates the introduction of channel memory from the
channel-estimation uncertainty introduced in the next layer.


\paragraph{Layer 3: finite channel estimation.}

Layer~3 retains the structured channel model but replaces exact parameter
knowledge with estimation from a finite pilot budget. Channel knowledge is
therefore converted from a free assumption into an acquired resource with
measurable estimation error and overhead.

An important distinction in this layer is between \emph{structural
identifiability} and \emph{operationally reliable identification}. A channel
may remain identifiable in principle while the available pilot budget produces
estimates that are too uncertain to support reliable matched detection.

The pilot sweep therefore measures when the information advantage of the
model-based receiver becomes practically usable rather than assuming that
channel knowledge is available without cost.


\paragraph{Layer 4: no supplied per-block matched posterior.}

Layer~4 removes the supplied per-block channel information required to
construct the matched model-based posterior.

The learned relay is not assumed to have no prior information about the
impairment family. It retains information acquired offline from its training
distribution. The evaluated blind classical methods, by contrast, must infer
useful structure from the received block without being supplied with the
per-block information required by the matched detector.

The comparison is therefore between amortized family-level learning and
per-block blind adaptation. This distinction is important because the learned
relay's training distribution constitutes prior information even when no
explicit channel estimate or matched posterior is supplied during inference.


\section{Relationship Between Identifiability and Arithmetic}
\label{sec:identifiability-arithmetic}

The four-layer sequence primarily locates the identifiability boundary. The
arithmetic boundary is orthogonal to that sequence.

Identifiability asks whether the information required to instantiate the
classical rule is supplied or can be estimated reliably. Arithmetic asks
whether that rule, once constructible, can be executed within the available
computational and decision-delay budget.

The two boundaries therefore define three broad operating regimes:

\begin{enumerate}
    \item the classical rule is constructible and affordable;

    \item the classical rule is constructible but becomes increasingly
    expensive as the problem dimension, channel memory, modulation order, or
    latency requirement changes;

    \item the matched classical rule is not operationally constructible from
    the information supplied at inference.
\end{enumerate}

The thesis does not assume that learned inference has complexity independent
of channel memory. For a fixed deployed architecture, its per-symbol inference
cost is fixed, but the context required by that architecture may increase as
the channel memory increases. The relevant comparison is therefore one of
scaling. Exact sequence-search complexity grows rapidly with channel memory
and modulation order, whereas the evaluated learned processing can exhibit
more favorable scaling with the context required by the deployed model.

Computational and latency analyses are consequently evaluated alongside the
relevant experiments rather than treated as an additional information layer.


\section{Research Hypotheses}
\label{sec:research-hypotheses}

The hypotheses below translate the thesis-level research question into
experimentally testable claims. H1--H4 concern the matched canonical control,
while H5 addresses the principal robustness question across Layers~2--4.


\paragraph{H1: Neural relays can improve upon AF in the low-SNR regime.}

Under the canonical matched benchmark, at least some evaluated learned relay
functions are expected to achieve lower BER than amplify-and-forward,
particularly at low SNR where forwarding amplified noise limits AF
performance.

H1 is deliberately restricted to comparison with AF. Evidence that a learned
relay improves upon AF does not imply that it improves upon matched
decode-and-forward.


\paragraph{H2: Learned relays do not reliably outperform matched symbol-wise
DF.}

Under the canonical matched, memoryless, and uncoded benchmark, no evaluated
learned relay is expected to achieve reliably lower BER than symbol-wise DF at
medium-to-high SNR. Learned relays may match DF within statistical uncertainty,
but are not expected to improve upon the matched classical decision rule.

H2 therefore establishes the expected control result when the information
required by the classical rule is fully available.


\paragraph{H3: Relay performance saturates beyond modest model capacity.}

Increasing neural model capacity beyond that required to represent the
canonical relay mapping is not expected to produce a corresponding BER
improvement.

H3 concerns performance saturation with model size. The experiments do not
assume in advance that any degradation observed at larger model sizes must be
caused by overfitting. Overfitting is therefore treated as a possible
mechanism only if the experimental evidence demonstrates a reproducible
degradation beyond statistical and initialization variation.


\paragraph{H4: Equal parameter budgets reduce but do not necessarily
eliminate architecture-dependent BER differences.}

When learned relays are compared under approximately equal parameter budgets
and input context, architectural differences are expected to narrow where the
underlying relay mapping is simple, particularly at high SNR. Differences may
remain at low SNR or when sequence structure becomes relevant.

Parameter normalization controls model size but does not isolate architecture
alone. Depth, width, state dimension, internal structure, and other
architecture-specific quantities may still differ between model families.


\paragraph{H5: Learned relaying becomes comparatively more useful as the
information required by matched classical processing is reduced.}

A learned relay trained across an impairment family is expected to become
comparatively more useful as the inference chain moves from exact channel
knowledge to finite estimation and finally to operation without the supplied
per-block information required to construct the matched posterior.

The expected advantage is not universal. Where the matched classical rule
remains both constructible and computationally affordable, it remains the
preferred reference. The hypothesis instead concerns the regimes in which the
information or resource assumptions required by that rule become restrictive.

For channels with memory, the comparison also depends on computational
scaling. Exact sequence-search complexity grows rapidly with channel memory
and modulation order, whereas learned inference may exhibit more favorable
scaling with the context required by the deployed architecture. H5 therefore
does not assume that learned complexity is independent of channel memory.


\section{Hypothesis-to-Experiment Mapping}
\label{sec:hypothesis-experiment-map}

Table~\ref{tbl:hypothesis-experiment-map} maps each hypothesis to the
experimental regime and primary evidence used to evaluate it.

\begin{table}[htbp]
\centering
\caption{Mapping between research hypotheses, experimental regimes, and
primary evidence.}
\label{tbl:hypothesis-experiment-map}
\begin{tabular}{@{}p{0.08\textwidth}p{0.19\textwidth}
                p{0.31\textwidth}p{0.32\textwidth}@{}}
\toprule
\textbf{Hyp.} &
\textbf{Regime} &
\textbf{Question} &
\textbf{Primary evidence} \\
\midrule

H1 &
Layer~1 &
Can learned relays improve upon AF? &
Canonical BER comparison \\

H2 &
Layer~1 &
Can learned relays reliably improve upon matched symbol-wise DF? &
Canonical BER comparison and analytical DF validation \\

H3 &
Layer~1 &
Does increasing neural capacity continue to improve BER? &
Model-size sweep and repeated training runs \\

H4 &
Layer~1 &
How much architecture-dependent behavior remains under comparable parameter
budgets? &
Normalized architecture comparison \\

H5 &
Layers~2--4 &
How does learned relaying compare as channel memory is introduced and matched
channel information becomes progressively less available? &
Known-ISI, finite-pilot, and posterior-free experiments \\

\bottomrule
\end{tabular}
\end{table}


\section{Decision Criteria}
\label{sec:hypothesis-decision-criteria}

The hypotheses are evaluated using the metrics appropriate to the question
being tested. BER versus SNR is the primary performance metric in the
canonical and unknown-channel comparisons. Model-size and architecture
hypotheses are interpreted using repeated training runs and the corresponding
statistical variation where available.

Computational comparisons use the operation-count or measured inference-cost
definitions specified in the corresponding experiment. Latency comparisons
distinguish computational cost from buffering and decision delay where block
or sequence processing is required.

A hypothesis is not interpreted as supported merely because one method
achieves a numerically lower BER at a single operating point. Conclusions are
interpreted with respect to the statistical uncertainty, SNR range,
information assumptions, and evaluation budget of the corresponding
experiment.

Likewise, a crossover between learned and classical methods is interpreted in
relation to the information supplied to each method. A comparison between a
matched detector with exact channel knowledge and a learned detector without
that information answers a different question from a comparison in which both
methods must operate from finite observations.


\section{Scope of the Hypotheses}
\label{sec:hypothesis-scope}

The hypotheses are empirical claims about the relay architectures, channel
models, modulation schemes, training procedures, and classical baselines
evaluated in this thesis. They are not universal claims about all neural or
classical communication systems.

In particular, failure of a learned relay to outperform matched DF in the
canonical benchmark does not imply that learned relaying is generally
ineffective. Conversely, an advantage under finite channel estimation or
operation without a supplied matched posterior does not imply that the same
learned relay should outperform a matched classical detector when the
information required by that detector is available.

References to the strongest classical baseline therefore mean the strongest
applicable classical baseline evaluated under the information assumptions of
the corresponding experiment. Claims of optimality are restricted to the
specified channel model, available information, and decision criterion.

The canonical benchmark, the principal unknown and mismatched-channel study,
and the supplementary coded and resource study also serve different purposes.
The canonical benchmark establishes the matched-model control. Layers~2--4
form the principal experimental contribution by progressively changing the
information conditions under which model-based processing operates. The coded
study is supplementary and tests the robustness of the canonical
interpretation under coding, throughput, and latency constraints. It does not
form an additional layer of the channel-information hierarchy.


\section{Chapter Synthesis}
\label{sec:research-objectives-synthesis}

The experimental program begins with the regime most favorable to classical
model-based processing and then progressively relaxes its preconditions.
Layer~1 establishes whether learning provides an advantage when the channel is
matched, memoryless, and known. Layer~2 introduces memory while retaining
exact channel knowledge. Layer~3 places a finite acquisition cost on that
knowledge. Layer~4 removes the supplied per-block information required to
construct the matched posterior.

These layers locate the information boundary of the comparison. Complexity
and latency analyses separately determine whether a classical rule that
remains statistically constructible can also be afforded operationally.

The hypotheses therefore do not seek a single winning relay architecture.
Together, they test the thesis proposition that the relative value of learned
and classical relaying depends on two conditions: whether the required model
information can be obtained reliably, and whether the resulting decision rule
can be executed within the available resource budget.

The chapters that follow evaluate these hypotheses first under the matched
canonical benchmark and then under progressively less favorable information
conditions. The resulting evidence is used to identify where classical
model-based processing remains the preferred reference and where learned
relaying becomes a useful complement.